Por definición, dado que no es posible observar directamente los estados ocultos de un Modelo Oculto de Márkov, la
probabilidad de obtener una secuencia específica de observaciones se calcula mediante la marginalización sobre el
espacio de todas las posibles secuencias de estados.\\
Utilizando la \textbf{ley de probabilidad total}, se establece que:

\[
p\!\left(w^{(1)} \cdots w^{(T)}\right) = \sum_{s^{(1)}} \cdots \sum_{s^{(T)}} p\!\left(w^{(1)} \cdots w^{(T)},\ s^{(1)} \cdots s^{(T)}\right)
\]

\medskip

Factorizamos la distribución conjunta aplicando la \textbf{regla de la cadena} de probabilidad, lo que permite
descomponer la secuencia en términos individuales:
\[
p\!\left(w^{(1)}\cdots w^{(T)},\ s^{(1)}\cdots s^{(T)}\right) = \prod_{i=1}^{T} p\!\left(w^{(i)},\, s^{(i)} \mid w^{(1:i-1)},\, s^{(1:i-1)}\right)
\]

\medskip

Para simplificar cada factor del producto, se aplican las dos hipótesis fundamentales que definen el
comportamiento de un HMM:

\begin{enumerate}
\item \textbf{Propiedad de Márkov:} Asumimos que el estado en el tiempo $i$ depende exclusivamente del estado
  inmediatamente anterior ($i-1$). Esto implica que el historial completo de estados previos no aporta información
  adicional para predecir el futuro inmediato:
  \[
  p\!\left(s^{(i)} \mid s^{(1)}, \ldots, s^{(i-1)}\right) = p\!\left(s^{(i)} \mid s^{(i-1)}\right) = A_{s^{(i-1)},\, s^{(i)}}
  \]

\item \textbf{Independencia de emisión:} Se supone que la observación actual depende únicamente del estado oculto
  presente, bajo esta premisa, ni las observaciones pasadas ni los estados previos influyen en la emisión actual
  si el estado actual es conocido:
  \[
  p\!\left(w^{(i)} \mid s^{(i)},\, w^{(1:i-1)},\, s^{(1:i-1)}\right) = p\!\left(w^{(i)} \mid s^{(i)}\right) = B_{s^{(i)},\, w^{(i)}}
  \]
\end{enumerate}

Al integrar estas suposiciones, se observa que cada componente de la factoría se reduce a la siguiente expresión:
\[
p\!\left(w^{(i)},\, s^{(i)} \mid w^{(1:i-1)},\, s^{(1:i-1)}\right) = p\!\left(w^{(i)} \mid s^{(i)}\right) \cdot p\!\left(s^{(i)} \mid s^{(i-1)}\right)
\]

\medskip

Finalmente, al sustituir este resultado en la sumatoria de la probabilidad total, se obtiene la fórmula general
requerida:
\[
p\!\left(w^{(1)} \cdots w^{(T)}\right) = \sum_{s^{(1)}} \cdots \sum_{s^{(T)}} \prod_{i=1}^{T} p\!\left(w^{(i)} \mid s^{(i)}\right) p\!\left(s^{(i)} \mid s^{(i-1)}\right)
\]

\vspace{1em}

Como mostró, debido a que los estados ocultos no son accesibles, es imperativo considerar
\textbf{todos los caminos posibles} a través del espaco de estados y sumar sus respectivas contribuciones.
Para cualquier camino específico $s^{(1)} \to s^{(2)} \to \cdots \to s^{(T)}$, la probabilidad de que dicho camino
genere la secuencia observada es el producto término a término de las probabilidades de transición y emisión:

\[
\underbrace{p(w^{(1)}\mid s^{(1)})}_{\text{emisión } 1} \cdot \underbrace{p(s^{(2)}\mid s^{(1)})}_{\text{transición}} \cdot \underbrace{p(w^{(2)}\mid s^{(2)})}_{\text{emisión } 2} \cdots \underbrace{p(s^{(T)}\mid s^{(T-1)})}_{\text{transición}} \cdot \underbrace{p(w^{(T)}\mid s^{(T)})}_{\text{emisión } T}
\]

\medskip

Es la estructura de producto derivada de las suposiciones de Márkov y de independencia lo que hace que este
modelo sea computacionalmente tratable. Sin ellas, el costo de cálculo crecería de forma exponencial respecto a
la longitud de la cadena. Esta misma arquitectura es la que aprovecha el \textbf{algoritmo de Viterbi}; la
diferencia fundamental radica en que, en lugar de sumar sobre todos los caminos, se selecciona el camino que
maximiza la probabilidad, permitiendo hallar la secuencia de estados más probable en un tiempo eficiente de
$O(T \cdot |S|^2)$.
